\chapter{System Design and Implementation}
\label{ch:system}

\section{Hardware Setup}

\subsection{Receivers}

Three ESP32-S3-WROOM-1 N16R8 development boards (16\,MB flash, 8\,MB PSRAM)
serve as receivers. Each runs the \texttt{esp\_csi} \texttt{csi\_recv\_router}
example \cite{espcsi} with two modifications that are necessary rather than
cosmetic, and which are kept as a patch file so they survive any re-clone of
the upstream repository:

\begin{enumerate}
  \item \textbf{WiFi power save disabled} (\texttt{WIFI\_PS\_NONE}). With
        power save active the station sleeps between beacons and misses the
        downlink frames that generate CSI, producing an irregular and
        substantially reduced rate.
  \item \textbf{A custom console UART.} The default USB-JTAG console routes
        CSI output to a USB endpoint that is not the one being read, and
        does not permit the baud rate to be configured. Only in custom UART
        mode can the console be set to 921600 baud.
\end{enumerate}

Receivers are identified by the serial number of their USB-UART bridge,
which is stable across reboots and re-plugs, so a given physical board maps
to the same logical identifier in every session without manual
configuration. Serial port, IP address, MAC address, and measured RSSI for
each board are recorded in the collection protocol document.

\subsection{Receiver placement}

The three receivers are positioned at different locations and heights around
the sensing area, with tripod feet marked on the floor so the geometry is
reproducible across sessions, and photographed before every session. The
placement and the router position are recorded in each session's metadata.

Geometry is treated as part of the experimental configuration rather than as
an incidental detail, because a change in receiver position changes the
multipath structure the model has learned; an unrecorded move between
sessions would appear in the results as an unexplained cross-session drop.

\begin{figure}[htbp]
  \centering
  \fbox{\parbox[c][0.30\textwidth][c]{0.85\textwidth}{\centering
    \texttt{room\_layout}\\[0.5em]
    Room plan: receiver positions and heights, router position,
    sensing area, tripod floor markings.\\
    To be inserted from the measured session layout.}}
  \caption{Receiver and router placement in the primary collection room.}
  \label{fig:layout}
\end{figure}

\subsection{Router configuration and the RF environment}

The transmitter is an ISP-supplied consumer router operating on the 2.4\,GHz
band at 20\,MHz bandwidth, WPA2-PSK. Two constraints of this unit are
documented because they affect the data rather than merely the convenience
of collecting it.

First, \textbf{the channel cannot be pinned}. The unit runs automatic
channel selection and its subscriber-level administration interface does not
expose a manual override, so the operating channel is wherever auto-select
has landed and a mid-session change cannot be excluded a priori. The
mitigation is detection rather than prevention: session quality assurance
checks every session for channel and bandwidth drift, and a session whose
frames straddle two channels fails and is discarded rather than entering the
dataset unnoticed.

Second, the channel that auto-select chose is \textbf{congested}. A
41-scan channel survey conducted at the collection site measured the
strongest competing access point on each channel:

\begin{table}[htbp]
  \centering
  \caption{Channel survey at the collection site (41 scans over
  approximately 7 minutes). The project's own access point is excluded from
  the interference figures --- it is signal, not noise.}
  \label{tab:channel-survey}
  \begin{tabular}{clp{0.45\textwidth}}
    \toprule
    Channel & Strongest other AP & Assessment \\
    \midrule
    1  & $-82$\,dBm & Usable fallback; furthest from microwave-oven emission \\
    5  & $-73$ to $-78$\,dBm & Overlaps channel 6 \\
    6  & $\mathbf{-45}$\,\textbf{dBm} & \textbf{In use} --- a neighbouring AP
         only 4\,dB below our own \\
    8  & $-93$\,dBm & Overlaps 6 and 11 \\
    11 & $-89$\,dBm, frequently undetected & \textbf{Preferred} \\
    \bottomrule
  \end{tabular}
\end{table}

Channel 6 places a neighbouring access point at $-45$\,dBm against the
project's own $-41$\,dBm --- near-equal-power co-channel contention. That
neighbour's traffic both consumes airtime and injects uncontrolled multipath
variation into every recorded window. Channel 11 would improve the
AP-interference floor by roughly 44\,dB.

The recommended remedy is a dedicated access point: it need not provide
internet connectivity, only transmit, and it resolves the channel pinning,
band steering, and competing-client problems simultaneously. This is
recorded as a known limitation of the collection environment and is
revisited in the failure analysis.

\section{Software Architecture}

The host software is a single Python package organised as a pipeline of
small, independently testable modules. The design rule throughout is
immutability: frozen dataclasses for configuration and records, pure
functions returning new arrays, and no shared mutable state between stages.

\begin{description}
  \item[Parser] Converts one \texttt{CSI\_DATA} line into a frame record
    (metadata plus a complex subcarrier vector). The byte layout is pinned
    by tests, including the imaginary-before-real ordering, so an upstream
    firmware change cannot silently alter the interpretation.

  \item[Collector] Reads the three serial ports concurrently on separate
    threads, timestamps each frame on arrival using the host clock, and
    writes per-receiver Parquet files. Host timestamping at arrival is what
    makes window-level alignment possible across three unsynchronised
    device clocks.

  \item[Traffic generator] Produces the sustained UDP downlink described in
    Section~\ref{sec:csi-acquisition}. It is exposed as a shared context
    manager used by every component that reads from live hardware, so that
    the requirement cannot be forgotten in one entry point while being
    honoured in others.

  \item[Storage] Defines the session layout: per-receiver Parquet files, a
    \texttt{metadata.json} sidecar, and a \texttt{labels.json} sidecar for
    scripted sessions. Sessions carry an explicit exclusion flag for
    rig-test captures, honoured by dataset assembly.

  \item[Preprocessing] Pure functions implementing the chain of
    Chapter~\ref{ch:methodology}, each unit-tested for correctness and for
    non-mutation of its input.

  \item[Dataset] Assembles sessions into tensors, applies labels with
    boundary trimming, and provides the four split functions together with
    the degenerate-coverage check.

  \item[Training] Model construction, the training loop, evaluation, the
    checkpoint contract, and the provenance record appended to the results
    CSV.

  \item[Realtime] Streaming inference (Section~\ref{sec:realtime}).

  \item[Dashboard] The demonstration interface (Section~\ref{sec:dashboard}).

  \item[Simulator] Generates synthetic CSI in the byte-exact firmware output
    format. This is what allows the entire downstream pipeline to be
    developed and regression-tested without hardware, and it is the reason
    the software was complete and green before the boards were flashed.
\end{description}

\subsection{Session preflight}

A ten-second live check is run immediately before every recording session,
verifying that all three receivers are present, that each is at rate, and
that all three are associated on the same channel. It exists because the
alternative --- discovering a dead board after a 25-minute recording --- is
expensive in a way that a ten-second check is not.

\section{Real-Time Inference}
\label{sec:realtime}

\subsection{Streaming design}

The live path maintains a ring buffer per receiver. On each tick, a 3\,s
window ending at the current time is sliced from each buffer, the three
windows are aligned and preprocessed through the \emph{same functions in the
same order} as the offline pipeline, the model is evaluated, and a
prediction is emitted. Ticks occur at the 1.5\,s hop.

Reusing the identical preprocessing functions --- rather than a
streaming-optimised reimplementation --- is a deliberate constraint. A
train/deploy preprocessing mismatch is a defect class that produces a
degraded live system while every offline test continues to pass, and it is
correspondingly difficult to diagnose.

\subsection{Measured latency}

Per-stage latency was benchmarked with three receivers at a 3\,s window and
1.5\,s hop, replaying a real capture. Against a hop budget of 1500\,ms, the
95th-percentile compute time for a complete tick is approximately 23\,ms ---
a margin of roughly 65$\times$. Within that, window construction (Hampel,
detrend, resample) accounts for about 85\,\%; model inference is
approximately 3\,ms; parsing and buffer ingest are approximately 0.02\,ms
per frame.

The engineering conclusion is that \textbf{end-to-end delay is structural,
not computational}: it is dominated by the 3\,s of signal that must
accumulate before a window exists at all. Optimising the code cannot
meaningfully reduce it; only shortening the window can, at a cost in
accuracy that the window-length ablation quantifies.

\subsection{A known and accepted offline/online gap}

The streaming engine scores several points below the offline pipeline on the
same recorded windows, with approximately 96\,\% window-for-window agreement
in predicted label. The cause is identified and is a genuine property of the
problem rather than an implementation defect.

The moving-mean detrend is a \textbf{centred} 101-tap kernel. Offline, every
interior sample of a window is detrended using 0.5\,s of data on each side
--- including 0.5\,s of data that lies in the sample's future. A live stream
cannot supply that. The implementation already takes 0.5\,s of lead-in
before each window so the leading edge is correct, but the trailing edge
cannot be corrected without waiting. The discrepancy is negligible in the
window interior and concentrated entirely in its final 0.5\,s.

Closing the gap is possible --- delay each prediction by 0.5\,s, which the
65$\times$ latency headroom easily affords --- but it trades additional
detection delay for a few points of accuracy. It is recorded here as a
deliberate design trade, and documented in the source, so that a future
maintainer does not spend effort "fixing" the window slicing in the
expectation of a large gain.

\subsection{Temporal smoothing}

Raw per-window predictions flicker: adjacent windows overlap by 50\,\% and a
borderline window can flip the displayed label several times during a single
continuous activity. A rolling majority vote over the last 5 predictions is
therefore applied, with a confidence threshold of 0.6.

The semantics of the threshold matter more than its value. A low-confidence
window \textbf{abstains} --- it is removed from the vote --- rather than
voting for a literal "unknown" label. The distinction is not academic: in
the earlier implementation "unknown" competed as a candidate in the majority
vote, so two unsure windows could out-vote three confident ones. That made
the smoother strictly worse than no smoothing at all: identical flicker and
several points lower accuracy. With abstention, "unknown" is displayed only
when nothing in the recent history is confident, which is the intended
meaning of the state.

Two honest observations about smoothing are reported rather than elided.
First, \textbf{smoothed accuracy is lower than raw accuracy}. Majority
voting suppresses independent noise, but recognition errors here are
temporally clustered --- a window the model gets wrong is usually adjacent
to other windows it gets wrong --- so the vote propagates errors as readily
as it corrects them. Second, the smoother's real contribution is
\textbf{display stability}: it substantially reduces the number of label
changes over a session, against a small number of genuine activity
transitions. The smoother therefore earns its place in the user interface,
not in the accuracy table, and the raw figure is the one quoted as the
system's recognition performance.

\subsection{Fall alert fast path}

A fall is transient: it occupies one or two windows. A 5-window majority
vote will therefore essentially never select \emph{falling}, which means
that raising the alert from the smoothed label structurally suppresses the
one output the system most needs to produce.

The fall alert consequently \textbf{bypasses the smoother}. Any single raw
window classified as \emph{falling} with confidence at or above the
threshold raises the alert immediately and independently of the vote. This
accepts a higher false-alarm rate on the safety-critical class in exchange
for not systematically missing the event, which is the correct direction for
the asymmetry of the two errors.

\section{Demo Dashboard}
\label{sec:dashboard}

The demonstration interface is a FastAPI application serving a single page
over a WebSocket connection. It displays live per-receiver CSI amplitude
heatmaps, the current raw and smoothed predictions with confidence, and a
latching fall-alert banner.

The page is \textbf{fully self-contained} --- no external stylesheets,
scripts, or fonts. This is a defence-day requirement rather than a
preference: the venue's network cannot be relied upon, and a dashboard that
fetches a resource from a CDN fails in exactly the circumstances where it
must work. Both the raw and smoothed labels are streamed, so the difference
discussed above is visible during the demonstration rather than hidden by
it.

\section{Quality Assurance}

\subsection{Session gates}

Every recorded session is checked before it is admitted to the dataset. A
receiver's capture fails if its mean rate falls below 90\,Hz, if its
sequence-number gap fraction exceeds 5\,\%, if the dominant modulation rate
is not sufficiently stable, if the channel or bandwidth drifts within the
session, or if the empirically detected null subcarriers do not match the
expected DC-and-guard-band set --- the last being a strong indicator that
something upstream is wrong with the parse or the radio configuration.

The check is run immediately after recording, so a bad session is discovered
the same day, while the room, the subject, and the operator are still
available to repeat it.

\subsection{Automated tests}

The package carries 247 automated tests. Their coverage is deliberately
weighted toward the properties that fail silently rather than loudly:

\begin{itemize}
  \item The parser's byte layout, including the imaginary-before-real
        ordering, pinned against captured firmware output.
  \item Non-mutation of inputs by every preprocessing function.
  \item That every live entry point drives the UDP traffic generator ---
        three components once shipped without it, and each would have
        recorded a silently worthless session.
  \item That the training path warns on degenerate splits and records the
        reason.
  \item That the results record stores the resolved compute device rather
        than the requested one.
  \item That low-confidence predictions abstain from the smoothing vote
        rather than voting.
\end{itemize}

Where a test guards a specific past defect, it was verified by mutation ---
reintroducing the defect and confirming that the test fails --- because a
regression test that passes under the bug it was written to catch provides
false assurance.

\subsection{Hardware-free development}

The simulator generates synthetic CSI in the exact byte format the firmware
emits, which allows every stage downstream of the serial port to be
developed and tested without boards. The practical consequence is that
hardware time --- the genuinely scarce resource, since it requires the room,
the rig, and the subjects simultaneously --- is spent on collecting data
rather than on debugging software.
